%% 03_introduction.tex — body only, no preamble

Accurate prediction of surgical case duration is a critical determinant of Operating Room (OR) efficiency, hospital throughput, and patient safety. When scheduled durations deviate substantially from actual case times, the consequences cascade through the entire surgical schedule: underestimates trigger overtime, staff fatigue, and delayed recovery care, while overestimates strand expensive theatre time and force the cancellation or postponement of elective procedures for patients already facing lengthy wait times \cite{oliveira2023systematic,park2025machine}. A systematic review of machine learning approaches to surgical duration prediction found that ML models achieve an average improvement of 25.7\% over traditional scheduling methods, underscoring the magnitude of the gap left by conventional heuristics \cite{park2025machine}.

Operating rooms are among the most expensive resources in any hospital, generating over 40\% of a hospital's revenue while accounting for nearly 30\% of all expenditures; OR time in the United States alone is estimated to cost between \$30 and \$90 billion annually, with each minute of OR time valued at approximately \$46 \cite{al2025comprehensive,park2025machine}. Even modest reductions in scheduling error therefore translate directly into substantial cost savings, improved staff productivity, and greater throughput for waiting patients \cite{wang2022more}.

Existing computational approaches to surgical duration prediction overwhelmingly rely on structured pre-operative Electronic Health Record (EHR) data — surgeon identifier, specialty, patient demographics, and procedure codes — while systematically discarding the free-text procedure description recorded at the time of booking \cite{kwong2025optimizing,park2025machine}. A systematic review of 11 studies published between 2019 and 2024 confirmed that this structured-only paradigm represents the current state of practice, with the best models reducing prediction error by up to 51\% yet still treating the procedure description as non-informative \cite{park2025machine}. Wang and Dexter further demonstrated that even modest reductions in mean absolute predictive error of OR times can yield meaningful increases in labour productivity when coupled with appropriate scheduling adjustments \cite{wang2022more}.

This omission is consequential: the procedure description is the single richest pre-operative signal, encoding the specific surgical plan with a granularity unavailable from any categorical code. Prior work demonstrated that incorporating procedure text into a multimodal prediction model reduces Mean Absolute Error (MAE) by 16\% relative to a structured-only baseline across 180{,}370 retrospective surgical cases, with Sentence-BERT and Bio-ClinicalBERT embeddings achieving the strongest performance \cite{noorchenarboo2026benchmarking}.

Incorporating procedure text effectively requires capable text encoders. Bidirectional Encoder Representations from Transformers (BERT) \cite{devlin2019bert} established state-of-the-art performance across a wide range of Natural Language Processing (NLP) tasks by pre-training deep bidirectional Transformer representations on large unlabelled corpora and fine-tuning with minimal task-specific layers. For sentence-level embedding tasks, Sentence-BERT (SBERT) \cite{reimers2019sentence} extends BERT with siamese and triplet network structures that produce semantically meaningful fixed-size sentence embeddings efficiently — reducing the time to find the most similar pair in a corpus of 10{,}000 sentences from 65 hours with raw BERT to approximately 5 seconds. For the clinical domain, Bio-ClinicalBERT \cite{alsentzer2019publicly} adapts BERT by pre-training on approximately two million clinical notes from the Medical Information Mart for Intensive Care (MIMIC-III) database, improving performance on clinical Named Entity Recognition (NER) and natural language inference tasks. However, both Bio-ClinicalBERT and SBERT carry approximately 110\,M and 22\,M parameters respectively, with disk footprints of 436\,MB and 91\,MB, making them impractical for deployment on the edge or mobile hardware that many hospital scheduling systems rely upon for point-of-care decision support.

Knowledge distillation offers a principled route to compact models that approximate the representational capacity of large teacher networks. Sanh et al.\ demonstrated with DistilBERT \cite{sanh2019distilbert} that a student model trained through a combination of distillation, masked language modelling, and cosine-embedding losses can reduce BERT's parameter count by 40\% and increase inference speed by 60\% while retaining 97\% of language understanding performance. Jiao et al.\ extended this framework with TinyBERT \cite{jiao2020tinybert}, proposing a two-stage distillation strategy — general distillation on a large corpus followed by task-specific fine-tuning — that produces a model 7.5$\times$ smaller and 9.4$\times$ faster than BERT-base on the General Language Understanding Evaluation (GLUE) benchmark. Despite these advances, no prior work has applied knowledge distillation to the surgical procedure text domain, nor evaluated the resulting model within a complete end-to-end surgical duration prediction pipeline that satisfies the latency and storage constraints of real OR scheduling systems.

This paper addresses this gap by proposing TinySurgicalBERT, a lightweight surgical-domain text encoder distilled from Bio-ClinicalBERT and compressed to 0.75\,MB — small enough for deployment on the resource-constrained hardware that OR scheduling systems commonly rely upon. The proposed model is evaluated against Bio-ClinicalBERT and Sentence-BERT across eight regression models on 180{,}370 retrospective surgical cases, demonstrating that domain-specific distillation achieves statistically equivalent predictive accuracy while requiring up to $580\times$ less storage. Together, these results establish that the deployment barrier of large clinical language models can be overcome without sacrificing prediction quality in the surgical scheduling domain. The remainder of this paper is organised as follows. Section~\ref{sec:related_work} reviews prior work on surgical duration prediction and clinical language modelling, identifying the gaps that motivate this work. Section~\ref{sec:methodology} details the data preprocessing, knowledge distillation, and experimental evaluation pipeline. Section~\ref{sec:results} reports results across all encoding and model combinations. Section~\ref{sec:conclusion} draws conclusions and outlines directions for future research.

%%
%% PDFs read: al2025comprehensive.pdf, oliveira2023systematic.pdf, kwong2025optimizing.pdf,
%%            noorchenarboo2026benchmarking.pdf, park2025machine.pdf, wang2022more.pdf,
%%            devlin2019bert.pdf, alsentzer2019publicly.pdf, reimers2019sentence.pdf,
%%            sanh2019distilbert.pdf, jiao2020tinybert.pdf
%% Citations used: oliveira2023systematic, park2025machine, al2025comprehensive, wang2022more,
%%                 kwong2025optimizing, noorchenarboo2026benchmarking, devlin2019bert,
%%                 reimers2019sentence, alsentzer2019publicly, sanh2019distilbert, jiao2020tinybert
